\section*{Cycle--Orbit Splitting: corrected status}

\noindent\textbf{Status (November 2026):} The previously asserted
cycle-rank-to-vertex-type-splitting implication is \emph{not proved}.
Its nonvacuous interpretation fails already for degree three and radius one.
The formal existential threshold statement is trivially true if its
threshold exceeds the maximal possible rank; that vacuity is not a locality
or type-splitting theorem.

For a nonempty finite simple graph $G$ of maximum degree $\Delta$, let
$B_{2R}(v)$ denote the induced radius-$2R$ ball and put
\[
\mathsf{CycComp}_R(G)
 = \max_{v\in V(G)}\dim_{\mathbb F_2} Z_1(B_{2R}(v)).
\]
For $\Delta\ge 2$ put
\[
M_{\Delta,R}=1+\Delta\sum_{j=0}^{2R-1}(\Delta-1)^j,
\qquad U_{\Delta,R}=1+\left\lfloor
 \frac{(\Delta-2)M_{\Delta,R}}{2}\right\rfloor .
\]
(The sum is empty when $R=0$.)

\noindent\textbf{Proposition (uniform ceiling).}
For every nonempty graph of maximum degree $\Delta\ge2$,
$\mathsf{CycComp}_R(G)\le U_{\Delta,R}$.
For $\Delta\le1$, every ball is acyclic and the cycle rank is zero.
\begin{proof}
A radius-$2R$ ball is connected and contains at most $M_{\Delta,R}$
vertices by breadth-first counting. With $n$ vertices, it has at most
$\lfloor\Delta n/2\rfloor$ edges. Its cycle-space dimension is therefore
$|E|-n+1\le 1+\lfloor(\Delta-2)n/2\rfloor \le U_{\Delta,R}$.
For $\Delta\le1$ the graph is a disjoint union of vertices and edges.
\end{proof}

\noindent\textbf{Proposition (nonvacuous splitting obstruction).}
Fix $\Delta=3$ and $R=1$. The Petersen graph, represented as the Kneser
graph $KG(5,2)$, has ten vertices, fifteen edges, degree three and diameter
two. Its radius-two ball at any vertex is the whole graph, with
$\mathbb F_2$ cycle rank $15-10+1=6=U_{3,1}$.
Permuting the five underlying symbols acts transitively on its ten
vertices, so \emph{every} isomorphism-invariant rooted vertex type agrees
at every vertex, including every ordinary $\mathrm{FO}^k$ type.
Thus no $t^*(k,3,1)\le6$ forces two different vertex types, whereas
$t^*(k,3,1)>6$ makes the antecedent impossible. This holds for every $k$.

\noindent\textbf{OPEN:} A genuinely nonvacuous type-separation principle
requires independently verified hypotheses that exclude homogeneous
examples, a precise logical signature (ordinary $\mathrm{FO}^k$ versus
counting logic), and an actual refinement operator tied to those semantics.
The former Cycle Fan Predicate / EF Cycle--Linkage Trap reference does
not discharge these obligations. In particular, no type-splitting result
is derived from a pigeonhole repetition statement alone.
